\subsection{Terminal triples (1,3,w) at a root (1,2,v)}\label{proof:Q5-C172}
% Common local aliases for the former repeated C168.1--C168.3 prefaces.
Use the original-measure calculus of
Section~\ref{sec:odd-common-calculus} and the complete root-pair
certificates of Section~\ref{sec:odd-common:fourth}.
For a root \((1,2,v)\), write \(m=v-2\), \(G_m=E+zD\),
\(\beta=\beta_{v-1}=\int D\,d\nu\),
\(\beta_v=\int E\,d\nu\),
\(P_+=h_2f_v\), \(P_-=\bar f_2W_v\),
\(\ell_+=\ell_a\), and \(\ell_-=\ell_s\).
The two vectors, their signed means and full constants are
\eqref{eq:odd-common:root-standard}--\eqref{eq:odd-common:root-losses}.
At \(v=3\), \((E,D)=(1,0)\) and \(\beta=0\).
All subsequent direction choices and packet estimates retain their
stated integer and parameter ranges.
Here $n=v-2$, and $D_r$ denotes the imaginary coordinate $D$. The ordinary-fourth bridge of Q5-C084, Section 2, has cost $11d$ for every positive outer length. Its exact terminal $(1,3)$ row is $\E[pP_{v,L}T^L(qf_3)]^4\le e^{-705d/128}$, with precisely the same $P_{v,L}$ as below. It supplies the high-domain comparison.
\paragraph{C172.0 The terminal family (1, 3, w)}

The construction retains the original tree, every arm, subdivision and equal-pair multiplicity, all independent positive connector lengths, the original edge count and arbitrary original label laws.

% Replace a pure model-definition block only; keep local root/terminal data.
Use the original finite-space model and arm messages of
Section~\ref{sec:odd-common-calculus}.
When the original \(K\)-notation is used, \(K=-B\) and
\(g=K\one=-u\); in either notation
\(\delta_j=\|B^j\one\|_2^2=\|K^j\one\|_2^2\).
All means, norms, adjoints and spectral integrals use the same
original probability measure.
The root residual is \((1,\allowbreak \allowbreak  2,\allowbreak \allowbreak  v)\), \(v\ge3\), and the terminal residual \((1,\allowbreak \allowbreak  3,\allowbreak \allowbreak  w)\), \(w\ge4\). The inner and outer lengths are independent positive integers. Fix
\[
F=qf_3h_w.
\]
Choose the root direction by v and L:
\[
P=P_{v,L}=\begin{cases}h_2f_v,&L=1,\ v\ge4,\\ \bar f_2W_v,&\text{all other cases},\end{cases}
\qquad H=pP\,T^LF.                                      \tag{C172.T179.1}
\]

For every listed original integer and legal kernel,
\[
\boxed{\begin{gathered}
EH^4\le e^{-705d/128},\qquad\\
I_q(T^MH)^3\le e^{-d/64},\quad I_q(Y)=E[q^{-1/3}|Y|^{8/3}].
\end{gathered}}\tag{C172.T179.2}
\]
The rate \(705/128=11/2+1/128\) is the ordinary fourth rate of the message containing the original unit arm and complete root pair. Restore genuine equal pairs and complete reversal by the common argument.

The root unit arm points outward, giving \(p\). The terminal unit arm points inward, its three-edge arm entirely outward, and its \(w\)-arm first inward then outward. Inner edges point from terminal to root and outer edges from root to receiver. On the entire low domain \(0<d\le1/4\), the more precise rates are
\[
\log EH^4\le\begin{cases}-421d/60,&L=1,\\-427d/60,&L=2,\\-337d/60,&L\ge3.\end{cases}            \tag{C172.T179.3}
\]
The common low-domain margin is \(337/60-705/128=209/1920>0\). The same root direction on the high domain compares with the exact Q5-C084 \(qf_3\) row.


\paragraph{C172.2 The actual terminal envelope and uniform low domain}

Until the final high-domain comparison, take \(0<d\le1/4\). Set
\[
m_0=w-3\ge1,\quad B_0=\beta_{w-1},\quad
\Gamma=d+\beta_w+B_0\ge(35/16)d,
\]
\[
R=-B_0u-B^3G_{m_0}(-B)u.                                \tag{C172.T179.8}
\]
Retain the actual constant \(B_0\) in the finite convolution of \(f_{w-1}\), then apply positive \(T\):
\[
h_w\le T S_{w-1}-B_0q,
\quad S_n=\sum_{j=0}^nK^j1,\ K=-B.
\]
Add the exact \(f_3\). Cancellation of the first three original \(K\)-powers gives
\[
f_3+h_w\le2-\Gamma+R.                                 \tag{C172.T179.9}
\]
At \(w=4\), \(\beta_3=\beta_4=d\), \(\Gamma=3d\), and \(R=-du-B^3u\); this is the exact short identity. Nonnegativity and AM--GM yield
\[
F=qf_3h_w\le q\exp(R-\Gamma).                          \tag{C172.T179.10}
\]

The original row and spectral bounds, retaining the \(G_2\) cost at \(m_0=2\), give
\[
|R|\le d+(6/5)\theta\sqrt d\le1/4+39/160=79/160<1/2.
\]
Thus \(e^R\le1+R+3R^2/5\). The integral Taylor coefficient increases for \(R\le1/2\), and \(e^{1/2}<33/20\) makes it less than \(3/5\); negative \(R\) is included. Use \(q\le2\), the original kernel bound two, and \(|B(uR)|\le\sqrt{dER^2}\):
\[
\begin{gathered}
T(qe^R)\le \bar f_2+BR+\delta,\\
\quad\delta=ER+E(uR)+\sqrt{dER^2}+(12/5)ER^2.
\end{gathered}\tag{C172.T179.11}
\]
The actual mixed mean remains signed.

\paragraph{C172.2.1 The case \(w=4\)}

Here \(ER+E(uR)=-\delta_2-d^2\), \(ER^2=d^3+\delta_4\). Put \(v=\delta_2/d\), \(J=ER^2/d\le d^2+\theta^2v\). Young gives \(\sqrt J\le1/8+2J\), hence
\[
\delta/d\le1/8-d+(22/5)d^2+v[-1+(22/5)\theta^2].
\]
The \(v\) coefficient is negative, since \(\theta^2<1/6\). Dropping it leaves a convex quadratic in \(d\), with endpoint values \(1/8,\allowbreak \allowbreak  3/20\) on \([0,\allowbreak \allowbreak  1/4]\). Hence \(\delta\le3d/20<7d/20\).

\paragraph{C172.2.2 The case \(w=5\) on the entire low domain}

Now \(B_0=d\), \(ER+E(uR)=-\delta_2-d^2+\delta_3\), and \(ER^2=d^3+\delta_4+\delta_5-2d\delta_3\). Thus \(J\le d^2+[\theta^2(1+\theta)-2\theta d]v\), with a positive middle coefficient. The same Young bound gives
\[
\begin{gathered}
\delta/d\le1/8-d+(22/5)d^2\\
+v[-\kappa+(22/5)\theta^2(1+\theta)-(44/5)\theta d].
\end{gathered}\tag{C172.T179.12}
\]
For \(d\le1/8\), drop the negative \(vd\) term and \(-d+22d^2/5\le0\), then bound a positive \(v\) coefficient by \(v\le\theta\). The result is \(156575/524288<3/10\). On \([1/8,\allowbreak \allowbreak  1/4]\), that coefficient decreases from \(-1593/81920<0\). Dropping it leaves the convex quadratic with endpoint values \(11/160,\allowbreak \allowbreak  3/20\). Thus the whole low domain has \(\delta\le7d/20\).

\paragraph{C172.2.3 Every w at least six}

In the original \(K\) notation let \(V_0=\sum_{j=4}^wK^j\one\), \(g=-u\), so \(R=B_0g-V_0\). Finite alternating moments give \(EV_0-E(gV_0)\ge \delta_2-2\delta_3\), \(E(gV_0)\ge\kappa \delta_3\), and the finite contracting \(G_{w-3}\) gives \(\|V_0\|^2\le \delta_4\). Therefore
\[
ER+E(uR)\le-(3/16)\delta_2-\kappa d^2,
\quad J\le d^2+[\theta^2-2\kappa^2\theta d]v.
\]
The coefficient \(\theta-2\kappa^2d\) is positive, so \(\delta_3\le\theta \delta_2\) is used in its valid direction. The same Young inequality gives
\[
\begin{gathered}
\delta/d\le1/8-\kappa d+(22/5)d^2\\
+v[-3/16+(22/5)\theta^2-(44/5)\kappa^2\theta d].
\end{gathered}\tag{C172.T179.13}
\]
The \(v\) coefficient stays positive for \(d\le1/4\); using \(2/5<\theta<5/12\), \(\kappa<3/5\), a lower bound is \(373/2000\). Set \(v=\theta\). The resulting convex quadratic has exact endpoints
\[
P(0)=28167/81920,
\qquad P(1/4)=1795141/5242880,
\]
Both are less than \(7/20\), whose numerator over \(5242880\) is \(1835008\). Convexity covers the whole domain and every \(w\ge6\).

Thus use the uniform actual first-row bound
\[
\boxed{D=7d/20\ge\max(\delta,0)}.                       \tag{C172.T179.14}
\]
This uniform certificate keeps the same original direction.

\paragraph{C172.3 Every finite signed moment and positive propagation}

For i at least one, let \(D_i=EB^iR.\) Finite alternating sums give
\[
\begin{gathered}
D_{2r-1}=(-1)^{r+1}[B_0\delta_r-H_{r+2}],\qquad\\
D_{2r}=(-1)^{r+1}\widetilde H_{r+2},
\end{gathered}\tag{C172.T179.15}
\]
Here \(H,\allowbreak \allowbreak  \widetilde H\) have respectively \(\lfloor m_0/2\rfloor,\allowbreak \allowbreak  \lceil m_0/2\rceil\) terms, starting at the indicated \(\delta_{r+2}\), and empty sums are zero. Hence \(0\le H,\allowbreak \allowbreak  \widetilde H\le \delta_{r+2}\), including \(m_0=1\). Consequently
\[
\begin{gathered}
D_1\le d^2,\quad0\le D_2\le \delta_3,\quad D_3\le \delta_4,\\
\quad D_4\le0,\quad(D_i)_+\le d/24\quad(i\ge4).
\end{gathered}\tag{C172.T179.16}
\]
For odd \(i\ge5\), the two summands have opposite signs, so the positive part is at most \(\max(B_0\delta_3,\allowbreak \allowbreak  \delta_5)\). Since \(B_0\le d\le1/4\), \(\theta^2<1/6\), the first is less than \(d/24\); the second is at most \(\theta^4d<d/36\). Positive even indices start with \(D_6\le \delta_5\).

Put \(\gamma_j=\beta_j-D_{j-2}.\) For every actual j at least three,
\[
\boxed{\gamma_j\ge d/2}.                              \tag{C172.T179.17}
\]
The precise short bounds are \(\gamma_3\ge3d/4\), \(\gamma_4\ge(1-\theta^2)d\), \(\gamma_5\ge(\kappa-\theta^3)d>49d/96\), and \(\gamma_j\ge53d/96\), \(j\ge6\).

Original positive propagation gives, for every L at least one,
\[
\begin{gathered}
T^LF\le e^{-\Gamma}Y_L,\\
\quad Y_L=\bar f_{L+1}+B^LR+D\bar f_{L-1}\\
+\sum_{i=1}^{L-1}D_i \bar f_{L-1-i}.
\end{gathered}\tag{C172.T179.18}
\]
\(Y_L\) dominates a strictly positive actual message. For L at least two, the finite convolution is
\[
Y_L=1+G_{L+1}(B)u+B^LR+D\bar f_{L-1}
-\sum_{j=3}^{L+1}\gamma_j \bar f_{L+1-j}.
\]
Every gamma is nonnegative. Keep the final original \(\bar f_0\) term and use \(\bar f_j\le2\) for the positive \(D\) term:
\[
\begin{gathered}
Y_L\le1+u+c_L+B Q_L(B)u,\\
\quad c_L=2D-\gamma_{L+1},\\
\quad Q_L=G_L+z^{L-1}r(z),\quad R=r(B)u.
\end{gathered}\tag{C172.T179.19}
\]
The exact mean is \(EBQ_L(B)u=-\beta_{L+2}+D_L=-\gamma_{L+2}\). At \(L=1\), retain \(Y_1=1+u+D+B(u+R)\), using \(D\) only once.


Combine with the complete root-pair envelope and the curvature payment of the original p. For L at least two this gives
\[
\begin{gathered}
\log EH^4\le-4\Gamma+4(2D-\gamma_{L+1}-\gamma_{L+2})\\
-4\ell_\sigma+(17/4)\|[Q_L-zG_n(\mp z)]u\|^2+(10/3)\beta^2.
\end{gathered}\tag{C172.T179.23}
\]
The positive root uses \(G_n(-z)\), the negative one \(G_n(z)\); every cross term belongs to the same actual kernel. At \(L=1\), replace the scalar bracket by \(D-d+D_1\) and \(Q_L\) by \(1+r\). Since \(\beta\le d\le1/4\), the curvature cost is at most \(5d/6\).

\paragraph{C172.5 The two short-connector joint-frequency payments}

Write \(r=-t+zs\). At \(w=4\), \(t=d,\allowbreak \allowbreak  s=x\). For \(w\ge5\), write \(G_{w-3}(-z)=U-zV\), with \(1-x\le U,\allowbreak \allowbreak  V\le1\), so
\[
t=B_0-x^2V,\qquad s=xU.
\]
The common range is
\[
-x^2\le t\le d\le1/4,\qquad0\le s\le x.                \tag{C172.T179.24}
\]
The coordinate \(t\) may be negative here. The actual means are \(D_1=\int t\), \(D_2=\int xs\).

\paragraph{C172.5.1 The first connector, including its actual mean}

For the negative root the frequency expression is
\[
\begin{gathered}
J_1^-=4t-4[E+(1+x)D_r]\\
+(17/4)[(1-t+xD_r)^2+x(s-E)^2].
\end{gathered}\tag{C172.T179.25}
\]
For \(n\ge2\), both root-coordinate derivatives are negative. The \(D_r\)-derivative is bounded by \(-4(1+x)+(17/2)x(1+x+x^2)\), increasing to \(-13051/65536<0\) at \(\theta\); the \(E\)-derivative is at most \(-4+17\theta/2=-35/64\). Thus take \(D_r=E=1-x\). Since \(s\le x<1-x\), take \(s=0\). Convexity in \(t\), and derivative at most \(4-(17/2)(3/4)<0\) on \([0,\allowbreak \allowbreak  1/4]\), reduce to \(t=0,\allowbreak \allowbreak -x^2\). At zero the expression is
\[
\mathcal B(x)=-15/4+(67/4)x-(35/4)x^2-(17/4)x^3+(17/4)x^4<3/2.
\]
This increases with derivative at least \(77/16\) for \(x\le1/2\), and endpoint margin \(-245759/4194304\) below \(3/2\). At \(t=-x^2\), it is
\[
-15/4+(67/4)x-(17/4)x^2+(17/4)x^3
\le-15/4+871/128+17/48=1309/384<7/2.
\]
This drops a negative quadratic term and uses \(\theta^3<1/12\). At \(n=1\), retain \((E,\allowbreak \allowbreak  D_r)=(1,\allowbreak \allowbreak  0)\). For \(t\ge0\), the maximum at zero is \(1/4+17x/4<2\); at \(t=-x^2\), it is at most \(1/4+9x^2/2+17x^4/4+17x/4<26/9<7/2\). Hence every negative-root first row costs at most \(7/2\).

For the actual \(L=1,\allowbreak \allowbreak  v\ge4\) positive root, replace the real coordinate by \(1-t-xD_r\), and root loss by \(E+(1-x)D_r\). The real coordinate is at least \(1-1/4-\theta>0\), so both root derivatives are negative. Take \(D_r=E=1-x\). Throughout the actual \(t\) range, \(0<1-t-x(1-x)\le1-x+2x^2\le1\), and \(4t\le1\). Set \(s=0\), obtaining
\[
1-969/256+(17/4)(1+\theta)=817/256<7/2.                 \tag{C172.T179.26}
\]
This covers precisely the actual first-row directions.

\paragraph{C172.5.2 The second connector, including both actual means}

Use the negative root, with \(Q2=1-xs+z(1-t).\) The joint expression is
\[
\begin{gathered}
J_2=4t+4xs-4[E+(1+x)D_r]\\
+(17/4)[(1+x(D_r-s))^2+x(1-t-E)^2].
\end{gathered}\tag{C172.T179.27}
\]
The \(D_r\)-derivative is at most \((1+x)(17x/2-4)<0\); the \(E\)-derivative is \(-4+(17/2)x(E+t-1)<0\). For \(n\ge2\), set \(D_r=E=1-x\); at one retain \((0,\allowbreak \allowbreak  1)\). The \(s\)-derivative is \(x[4-(17/2)(1+x(D_r-s))]<0\). Indeed even enlarging to \(s\le1\) leaves the real coordinate at least \(1-x>8/17\).

For nonnegative \(t\), enlarge to \([0,\allowbreak \allowbreak  1]\), keeping the \(E\)-derivative negative. Convexity reduces to endpoints; one exceeds zero by \(4+17x(2E-1)/4>0\). At \(s=0,\allowbreak \allowbreak  t=1\), the \(n=1\) value is \(17(1+x)/4\le765/128<6\), and later values are \(4+\mathcal B(x)<11/2\).

For negative \(t\), only the endpoint \(-x^2\) remains. The real coordinate is at most \(1+x\), and the imaginary coefficient \(x+x^2\) is at most one; this also bounds \(n=1\). Root loss is at least one, so
\[
\begin{gathered}
J_2\le-4x^2-4+(17/4)[(1+x)^2+x]\\
=1/4+(51/4)x+x^2/4\le22409/4096<6.
\end{gathered}\tag{C172.T179.28}
\]
Thus the full second-row cost is at most six, including negative t.

\paragraph{C172.6 Every finite connector of length at least three}

The root is negative. Write \(G_L=U+zV\), with both coordinates in \([1-x,\allowbreak \allowbreak  1]\) and \(U\le1-x+x^2\). The joint base satisfies
\[
-4[E+(1+x)D_r]+(17/4)[(U+xD_r)^2+x(V-E)^2]\le1/4.       \tag{C172.T179.29}
\]
At \(n=1\), its square is at most \((1-x+x^2)^2+x^3\le1\), since \(2-3x+x^2-x^3>0\). At \(n\ge2\), both root derivatives are negative. Set \(D_r=E=1-x\); the real coordinate is at most one, imaginary coefficient has absolute value at most \(x\), giving \(-15/4+4x+4x^2+17x^3/4\le-7/32<1/4\).

Uniformly for every finite n at least one and L at least three,
\[
\begin{gathered}
|G_L-zG_n(z)|^2\le(1+x)^2+x^3\\
\le66997/32768<67712/32768=(23/16)^2.
\end{gathered}\tag{C172.T179.30}
\]
The actual \(r\) is finite. Retaining every \(G_{m_0}\) bound \(6/5\) and \(\sqrt\theta<13/20\),
\[
|r(z)|\le d+(6/5)x^{3/2}
\le1/4+(6/5)\theta^{3/2}<907/1600<3/5.
\]
The perturbation modulus is at most \(3x/5\), for every finite \(L\ge3\). Its additional square is at most \(69x/40+9x^2/25\). After integration, the norm cost is
\[
\begin{gathered}
(17/4)[(69/40)\theta+(9/25)\theta^2]d\\
=330837d/102400<13d/4.
\end{gathered}\tag{C172.T179.31}
\]
This one certificate covers lengths three, four and every finite longer length.

\paragraph{C172.7 Continuous low-domain payments and the same-direction high domain}

For \(L=1\), the exact common payment is \(-4\Gamma+4D-4d+\int J_1+10\beta^2/3\). Use \(\Gamma\ge35d/16,\allowbreak \allowbreak  D=7d/20,\allowbreak \allowbreak  J_1\le7/2\), and curvature cost \(5d/6\):
\[
\log EH^4\le[-35/4+7/5-4+7/2+5/6]d=-421d/60.
\]
For \(L=2\), the adjacent gamma sum is \(2d-D_1-D_2\), so
\[
\log EH^4\le[-35/4+14/5-8+6+5/6]d=-427d/60.
\]
For every \(L\ge3\), the adjacent gamma sum is at least \(d\), with base cost \(d/4\) and additional norm cost \(13d/4\):
\[
\log EH^4\le[-35/4+14/5-4+1/4+13/4+5/6]d=-337d/60.
\]
Every low-domain row therefore has ordinary fourth loss greater than \(705/128\).

For \(d\ge1/4\), keep the same \(F,\allowbreak \allowbreak  P\). Since \(w-1\ge3\),
\[
h_w\le2Ef_{w-1}\le2(1-2d)\le1.
\]
Positivity gives \(0\le H\le pPT^L(qf_3)\). The root is exactly the Q5-C084 direction: positive at \(L=1,\allowbreak \allowbreak  v\ge4\), negative otherwise. Its same-direction fourth rate is \(705/128\). At \(d=1/4\), both analytic domains include the endpoint; at \(d=0\), original leaf integration gives every message one even if \(B\ne0\).

For every independent positive outer length use the Q5-C084 ordinary fourth bridge:
\[
I_q(T^MH)^3\le e^{11d}(EH^4)^2
\le\exp\{[11-2(705/128)]d\}=e^{-d/64}.
\]
The input is the actual ordinary fourth moment of the complete message.
